\chapter{Operacje kohomologiczne: naturalność, iloczyny i relacje}
\label{ch:operacje-v8}

\section{Dlaczego potrzebujemy operacji, a nie tylko pierścienia kohomologii?}
Iloczyn kubkowy porównuje dwie klasy tej samej przestrzeni. Kwadraty
Steenroda są silniejszym narzędziem: każda operacja
$\operatorname{Sq}^i:H^n(X;\F_2)\to H^{n+i}(X;\F_2)$
jest naturalna względem $X$ i zależy od jednej klasy.
Jeżeli dwie przestrzenie mają izomorficzne pierścienie kohomologii,
nie wynika jeszcze, że identyfikacja zachowuje wszystkie operacje
kohomologiczne. Szczególnie wyrazisty przykład otrzymamy po zawieszeniu
$\C P^2$, kiedy wszystkie iloczyny dodatnich stopni znikną, ale
$\operatorname{Sq}^2$ pozostanie nietrywialne.

W poprzednim rozdziale skonstruowaliśmy operacje za pomocą wyższych
diagonalnych. Należy przy tym oddzielić dwa twierdzenia:
acykliczność kompleksu modeli pozwala rozwiązać relacje homotopii
diagonalnych, ale \emph{nie wymusza} wybranego unormowania ani
wszystkich zgodności iterowanych diagonalnych. Stosując klasyczne
\emph{unormowane operacje Steenroda}, wykorzystujemy także te dodatkowe
zgodności. Właśnie one są potrzebne do wzoru Cartana i relacji Adema.

\section{Operacje naturalne i kobrzeg Bocksteina}
\begin{definicja}[Operacja kohomologiczna]
Niech $G,H$ będą grupami abelowymi, a $n,m\ge0$.
\emph{Naturalną operacją kohomologiczną} stopnia $(n,m)$ nazywamy
rodzinę funkcji
$\theta_X:H^n(X;G)\to H^m(X;H)$, indeksowaną
przestrzeniami topologicznymi $X$, taką że dla każdego ciągłego
$f:X\to Y$ i $u\in H^n(Y;G)$ zachodzi
$f^*\theta_Y(u)=\theta_X(f^*u)$.
Nie wymagamy addytywności, o ile nie jest jawnie wskazana.
Jeżeli $m=n+i$, mówimy, że operacja \emph{podwyższa stopień o $i$}.
\end{definicja}

\begin{przyklad}[Dlaczego mnożenie nie jest zawsze addytywne]
Przyporządkowanie $u\mapsto u\smile u$ jest naturalne. Dla klas
całkowitych stopnia parzystego nie musi być addytywne:
$(u+v)^2=u^2+2uv+v^2$. Przy współczynnikach w $\F_2$ składniki mieszane
znoszą się, ponieważ iloczyn jest przemienny w charakterystyce $2$.
Natomiast addytywność wszystkich $\operatorname{Sq}^i$
nie wynika z samej naturalności: jest osobną własnością
standardowej konstrukcji operacji Steenroda.
\end{przyklad}

\begin{definicja}[Kobrzeg Bocksteina]
Dla krótkiego ciągu dokładnego grup abelowych
$0\to A\xrightarrow{\iota}B\xrightarrow{q}C\to0$
\emph{homomorfizmem Bocksteina} nazywamy łączący homomorfizm
$\beta:H^n(X;C)\to H^{n+1}(X;A)$
w długim ciągu kohomologii współczynników.
Dla $0\to\F_2\xrightarrow{\,\cdot2\,}\Z/4\Z\to\F_2\to0$
piszemy $\beta_2:H^n(X;\F_2)\to H^{n+1}(X;\F_2)$.
\end{definicja}

\begin{stwierdzenie}[Wzór wykonawczy dla Bocksteina]
\label{prop:bockstein-v8}
Niech $a\in Z^n(X;\F_2)$ i niech $\widetilde a\in C^n(X;\Z/4\Z)$
będzie dowolnym podniesieniem kołańcucha $a$. Wtedy
$\delta\widetilde a$ przyjmuje wartości w podgrupie
$2(\Z/4\Z)$, a klasa kołańcucha
$(\delta\widetilde a)/2$ zidentyfikowanego z kołańcuchem
$\F_2$ jest równa $\beta_2[a]$.
\end{stwierdzenie}
\begin{proof}
Ponieważ $q(\delta\widetilde a)=\delta(q\widetilde a)=\delta a=0$,
wartości $\delta\widetilde a$ należą do jądra $q=2(\Z/4\Z)$.
Wybierzmy $b\in C^{n+1}(X;\F_2)$ spełniający
$\iota b=\delta\widetilde a$. Z równości
$0=\delta^2\widetilde a=\iota\delta b$ i różnowartościowości
$\iota$ dostajemy $\delta b=0$.
Sprawdźmy oba rodzaje wyborów. Dwa podniesienia tego samego $a$
różnią się o $\iota c$, gdzie $c\in C^n(X;\F_2)$;
wówczas odpowiadający im kołańcuch $b$ zmienia się o $\delta c$.
Jeśli natomiast zmienimy reprezentanta na $a'=a+\delta v$,
to wybieramy podniesienie $\widetilde v\in C^{n-1}(X;\Z/4)$
i kładziemy $\widetilde a'=\widetilde a+\delta\widetilde v$.
Z równości $\delta^2\widetilde v=0$ wynika, że otrzymujemy
ten sam $b$. Dowolne inne podniesienie $a'$ różni się
od niego o element obrazu $\iota$, więc jego wynik
reprezentuje tę samą klasę kohomologii. Jest to
kołańcuchowa konstrukcja homomorfizmu łączącego.
Ponieważ grupy łańcuchów singularnych są wolne,
krótki ciąg współczynników daje krótki ciąg kompleksów
kołańcuchów i opisany argument obowiązuje dla dowolnego $X$.
\end{proof}

\begin{przyklad}[Bockstein liczony na komórkach $\R P^2$]
\label{prz:bockstein-rp2-komorki-v8}
Rozważmy $\R P^2$ z komórkami
$e^0,e^1,e^2$, po jednej w każdym wymiarze.
Komórkowy homomorfizm brzegowy
\[
 \partial_2:C_2(\R P^2;\Z)\longrightarrow C_1(\R P^2;\Z)
\]
jest mnożeniem przez $2$, a $\partial_1=0$.
Dla współczynników $\F_2$ wszystkie koróżniczki
komórkowe są zerowe; stąd
$H^1(\R P^2;\F_2)\cong H^2(\R P^2;\F_2)
\cong\F_2$.
Niech $t$ będzie klasą kołańcucha $a$, który
na $e^1$ ma wartość $1\in\F_2$.
Podnieśmy go do $\widetilde a\in C^1(\R P^2;\Z/4)$,
przyjmując $\widetilde a(e^1)=1$ modulo $4$.
Wówczas
\[
 (\delta\widetilde a)(e^2)
 =\widetilde a(\partial_2e^2)=2\pmod4.
\]
Z definicji Bocksteina dzielimy ten \emph{parzysty}
kołańcuch przez $2$ i redukujemy wynik modulo $2$.
Otrzymujemy kołańcuch przyjmujący wartość $1$
na $e^2$, czyli niezerowy generator stopnia $2$.
Z wcześniej obliczonego pierścienia
$H^*(\R P^2;\F_2)\cong\F_2[t]/(t^3)$
wiemy, że jego niezerową klasą stopnia $2$
jest $t^2$; sam kompleks komórkowy oblicza
grupy, ale nie identyfikuje jeszcze ich iloczynu.
Stąd
\[
 \beta_2(t)=t^2\ne0
 \quad\text{w }H^2(\R P^2;\F_2).
\]
To niezależna kontrola tożsamości
$\operatorname{Sq}^1(t)=t^2$.
Nie przeczy ona $\beta_2^2=0$:
$H^3(\R P^2;\F_2)=0$, więc następne
zastosowanie homomorfizmu musi dać zero.

Dla $\R P^3$ ten sam rachunek ma istotną różnicę.
Komórkowe $\partial_3=0$, natomiast $\partial_2=2$,
więc $\beta_2(t)=t^2\ne0$ nadal, lecz
$\beta_2(t^2)=0$ już na poziomie
koróżniczki podniesionego kołańcucha
z komórki $e^2$ do $e^3$.
\end{przyklad}

\begin{figure}[!htbp]
\centering
\begin{tikzpicture}[>=Stealth,font=\small]
 \node[draw=manifoldblue!65!black,rounded corners=3pt,
       fill=manifoldblue!8,minimum width=3cm,
       minimum height=.95cm,align=center] (a) at (0,0)
       {$C^1(\R P^2;\Z/4)$};
 \node[draw=manifoldblue!65!black,rounded corners=3pt,
       fill=manifoldblue!8,minimum width=3cm,
       minimum height=.95cm,align=center] (b) at (5,0)
       {$C^2(\R P^2;\Z/4)$};
 \node[draw=manifoldgreen!70!black,rounded corners=3pt,
       fill=manifoldgreen!8,minimum width=3cm,
       minimum height=.95cm,align=center] (c) at (0,-1.9)
       {$C^1(\R P^2;\F_2)$};
 \node[draw=manifoldgreen!70!black,rounded corners=3pt,
       fill=manifoldgreen!8,minimum width=3cm,
       minimum height=.95cm,align=center] (d) at (5,-1.9)
       {$C^2(\R P^2;\F_2)$};
 \draw[->,thick] (a)--(b)
       node[midway,above] {$\delta^1=\cdot2$};
 \draw[->,thick] (c)--(d)
       node[midway,below] {$\delta^1=0$};
 \draw[->,thick] (a)--(c) node[midway,left] {$\rho_2$};
 \draw[->,thick] (b)--(d) node[midway,right] {$\rho_2$};
 \node[align=center] at (2.5,-3.05)
       {$\widetilde a(e^1)=1
       \ \Longrightarrow\ (\delta\widetilde a)(e^2)=2
       \ \Longrightarrow\ \beta_2(t)=t^2$};
\end{tikzpicture}
\caption{Kołańcuchowy rachunek Bocksteina na $\R P^2$.
Po podniesieniu kocyklu do $\Z/4$ koróżniczka
jest podzielna przez $2$, a jej iloraz daje niezerowy
generator $H^2(\R P^2;\F_2)$.}
\label{fig:bockstein-rp2-v8}
\end{figure}

\begin{twierdzenie}[Identyfikacja pierwszego kwadratu]
\label{thm:sq1-bock-v8}
Dla każdej przestrzeni topologicznej $X$ i każdego $n\ge0$
zachodzi $\operatorname{Sq}^1=\beta_2$ jako homomorfizm
$H^n(X;\F_2)\to H^{n+1}(X;\F_2)$.
W szczególności $\operatorname{Sq}^1\operatorname{Sq}^1=0$.
\end{twierdzenie}
\begin{proof}
Pierwsza równość jest \emph{standardową tożsamością unormowanej
konstrukcji $\smile_i$}. Jej kołańcuchową weryfikację otrzymuje się,
porównując wzór z poprzedniego rozdziału
$\operatorname{Sq}^1[a]=[a\smile_{n-1}a]$
z kołańcuchem $(\delta\widetilde a)/2$ z
Stwierdzenia~\ref{prop:bockstein-v8}. Porównanie wymaga jawnego
unormowanego modelu $\smile_i$; sama relacja
$\partial\Delta_i+\Delta_i\partial=(1+T)\Delta_{i-1}$
nie jest do niego wystarczająca. Drugą równość można udowodnić bez tego porównania.
Niech $\rho_2:H^n(X;\Z)\to H^n(X;\F_2)$
oznacza redukcję modulo $2$, a
$\beta_{\Z}:H^n(X;\F_2)\to H^{n+1}(X;\Z)$
homomorfizmem łączącym ciągu
$0\to\Z\xrightarrow{\cdot2}\Z\to\F_2\to0$.
Wprost z kołańcuchowej definicji obydwu homomorfizmów
łączących mamy $\beta_2=\rho_2\circ\beta_{\Z}$.
Dokładność długiego ciągu daje
$\beta_{\Z}\circ\rho_2=0$. Dlatego
$\beta_2^2=\rho_2\beta_{\Z}\rho_2\beta_{\Z}=0$.
\end{proof}

\begin{uwaga}[Zakres poprzedniego dowodu]
W identyfikacji $\operatorname{Sq}^1=\beta_2$ wykorzystujemy
znany fakt o \emph{konkretnych} unormowanych diagonalnych, a nie
przypisujemy go abstrakcyjnemu istnieniu dowolnej rodziny homotopii.
W pozostałych obliczeniach niezależność od modelu jest rozumiana
na poziomie operacji kohomologicznych, nie równości kołańcuchów.
\end{uwaga}

\section{Całkowita operacja i rachunki z użyciem wzoru Cartana}
\begin{definicja}[Całkowity kwadrat Steenroda]
Dla przestrzeni $X$ i klasy jednorodnej $x\in H^n(X;\F_2)$
przyjmujemy $\operatorname{Sq}(x)=\sum_{i=0}^n\operatorname{Sq}^i(x)$.
Przedłużamy ten wzór addytywnie na $H^*(X;\F_2)$.
Korzystamy z aksjomatów $\operatorname{Sq}^0=\id$,
$\operatorname{Sq}^i(x)=0$ dla $i>\deg x$ i addytywności
unormowanych operacji. Suma jest skończona dla każdej
klasy jednorodnej.
Jest to odwzorowanie grup addytywnych; składowa stopnia $n+i$
operacji $\operatorname{Sq}$ ma stopień $i$.
\end{definicja}

\begin{twierdzenie}[Postać multiplikatywna wzoru Cartana]
\label{thm:cartan-v8}
Dla dowolnej przestrzeni $X$ całkowity kwadrat jest
homomorfizmem pierścieni z jedynką:
$\operatorname{Sq}(x\smile y)=\operatorname{Sq}(x)\smile\operatorname{Sq}(y)$.
Ponadto $\operatorname{Sq}(1)=1$, ponieważ
$\operatorname{Sq}^0(1)=1$ i wszystkie
$\operatorname{Sq}^i(1)$ dla $i>0$ są zerowe.
Równość składowych stopnia $\deg x+\deg y+k$
dla klas jednorodnych $x,y$ jest
wzorem Cartana z Twierdzenia~\ref{thm:steenrod-cartan-v7}.
\end{twierdzenie}
\begin{proof}
Dokładnie rozdzielamy krok formalny od geometrycznego.
Korzystamy z następującego standardowego twierdzenia
pomocniczego o \emph{zewnętrznym kwadracie}: istnieje
naturalna konstrukcja klas ekwiwariantnych, multiplikatywna
względem iloczynu kubkowego, której ograniczenie do
przekątnej daje unormowane kwadraty Steenroda.
Nie dowodzimy tu tego twierdzenia ekwiwariantnego;
pokazujemy natomiast szczegółowo, jak wynika z niego Cartan.
W klasycznej konstrukcji wybiera się wolną, ściągalną
przestrzeń $E\Sigma_2$, gdzie
$\Sigma_2=\{1,\tau\}$ działa przez zamianę dwóch czynników.
Oznaczamy $B\Sigma_2:=E\Sigma_2/\Sigma_2$;
ma ona typ homotopijny $\R P^\infty$.
Następnie tworzy się przestrzeń orbit
$E\Sigma_2\times_{\Sigma_2}(X\times X)$.
Po ograniczeniu do diagonalnego odwzorowania
$X\to X\times X$ otrzymuje się klasę w
$H^{2n}(B\Sigma_2\times X;\F_2)$ o postaci
$\sum_{i=0}^n t^{n-i}\otimes\operatorname{Sq}^i(x)$,
przy czym $t\in H^1(B\Sigma_2;\F_2)$ jest generatorem.
Naturalność i formuła Künnetha nad ciałem zapewniają jednoznaczność
współczynników tej ekspansji.

Mnożenie zewnętrznych potęg jest zgodne z iloczynem klas,
po czym ograniczenie do przekątnej zachowuje iloczyny.
Dla $x,y$ daje to iloczyn wielomianów
$\bigl(\sum_i t^{n-i}\otimes\operatorname{Sq}^i(x)\bigr)
 \bigl(\sum_j t^{m-j}\otimes\operatorname{Sq}^j(y)\bigr)$.
Porównanie współczynników przy $t^{n+m-k}$
pozwala otrzymać żądany wzór.
\end{proof}
\begin{uwaga}[O rzeczywistym zakresie argumentu]
W dowodzie użyto \emph{twierdzenia o zewnętrznych potęgach}:
istnienia klasy i jej multiplikatywności, a także jej
identyfikacji z wcześniejszymi kołańcuchowymi kwadratami.
Te trzy własności wymagają technicznego dowodu ekwiwariantnego
na poziomie rezolwenty grupy $\Sigma_2$; opisaliśmy krok rachunkowy,
a nie pełny niezależny dowód istnienia potęgi. Daje to przejrzyste
wyprowadzenie Cartana \emph{z jawnie wskazanej konstrukcji},
bez ukrywania potrzebnego faktu pomocniczego.
\end{uwaga}

\begin{przyklad}[Rachunek na iloczynie rzeczywistych przestrzeni rzutowych]
Niech $X=\R P^{a}\times\R P^{b}$ i niech
$u,v\in H^1(X;\F_2)$ będą klasami z odpowiednich czynników.
Dla $0\le i\le a$, $0\le j\le b$ wzór Cartana daje
$\operatorname{Sq}^{k}(u^iv^j)=
\sum_{r+s=k}\binom ir\binom js u^{i+r}v^{j+s}$,
przy czym jednomiany przekraczające relacje
$u^{a+1}=v^{b+1}=0$ przyjmujemy za zero.
W szczególności dla $i=j=1$:
$\operatorname{Sq}^1(uv)=u^2v+uv^2$,
$\operatorname{Sq}^2(uv)=u^2v^2$.
\end{przyklad}

\section{Relacje Adema: co rzeczywiście wynika z przykładów?}
\begin{twierdzenie}[Relacje Adema w charakterystyce dwa]
\label{thm:adem-v8}
Dla liczb całkowitych $0<a<2b$ i każdej przestrzeni $X$
na $H^*(X;\F_2)$ zachodzi
\[
\operatorname{Sq}^a\operatorname{Sq}^b
 =\sum_{j=0}^{\lfloor a/2\rfloor}
 \binom{b-j-1}{a-2j}
 \operatorname{Sq}^{a+b-j}\operatorname{Sq}^{j},
\]
gdzie dwumiany redukujemy modulo $2$.
\end{twierdzenie}
\begin{uwaga}[Status dowodu]
Jest to twierdzenie klasyczne. Pełny dowód kołańcuchowy
wymaga porównania dwóch iterowanych operacji z działaniem
większej grupy symetrii (np. w modelu rozszerzonych potęg),
a więc danych nieobecnych w samej relacji dla $\smile_i$.
Nie przyjmujemy rachunku na jednym pierścieniu
$H^*(\R P^\infty;\F_2)$ za dowód \emph{uniwersalny}.
Natomiast poniżej rozpisujemy konsekwencje i sprawdzamy ich
zgodność z wykonawczymi formułami.
\end{uwaga}
\begin{przyklad}[Dwie pierwsze relacje]
Dla $(a,b)=(1,1)$ jedyny dwumian jest
$\binom{0}{1}=0$, stąd
$\operatorname{Sq}^1\operatorname{Sq}^1=0$.
Jest to także konsekwencja tożsamości
$\operatorname{Sq}^1=\beta_2$.
Dla $(a,b)=(2,2)$ mamy
$\binom{1}{2}=0$, $\binom{0}{0}=1$, więc
$\operatorname{Sq}^2\operatorname{Sq}^2=
\operatorname{Sq}^3\operatorname{Sq}^1$.
Na klasie $t^m\in H^m(\R P^\infty;\F_2)$
obie strony działają jako mnożenie przez
$\binom m2\binom{m+2}2$ oraz
$m\binom{m+1}3$, liczone modulo $2$.
Zgodność współczynników można sprawdzić z twierdzenia Lucasa,
ale sam ten rachunek nie stanowi jeszcze dowodu
równości wszystkich naturalnych operacji.
\end{przyklad}

\begin{przyklad}[Sprawdzenie relacji $\operatorname{Sq}^2\operatorname{Sq}^2
=\operatorname{Sq}^3\operatorname{Sq}^1$ w niskich stopniach]
\label{prz:adem-rp-niskie-v8}
Dla $t^m\in H^m(\R P^\infty;\F_2)$
oba wyrażenia mają stopień $m+4$.
Ze wzoru $\operatorname{Sq}^j(t^m)=\binom mj t^{m+j}$
otrzymujemy
\[
 \operatorname{Sq}^2\operatorname{Sq}^2(t^m)
  =\binom m2\binom{m+2}2\,t^{m+4},
 \qquad
 \operatorname{Sq}^3\operatorname{Sq}^1(t^m)
  =m\binom{m+1}3\,t^{m+4}.
\]
Współczynniki można sprawdzić \emph{bez używania}
relacji Adema ani wcześniejszych twierdzeń o
dwumianach modulo $2$.
W pierścieniu $\F_2[z]$ zachodzi
$(1+z)^2=1+z^2$ i
$(1+z)^4=1+z^4$.
Niech $m=4q+r$, gdzie $0\le r<4$.
Wtedy
\[
 (1+z)^m=(1+z^4)^q(1+z)^r
 \quad\text{w }\F_2[z].
\]
Współczynnik przy $z^2$ jest zatem równy
$\binom r2\bmod2$; jest on niezerowy
dokładnie dla $r=2,3$.
Natomiast $m+2$ ma resztę $0,1$ modulo $4$
wtedy, gdy $m$ ma resztę $2,3$.
Dlatego $\binom m2$ i $\binom{m+2}2$
\emph{nie mogą} być jednocześnie nieparzyste.

Dla prawej strony, jeśli $m$ jest parzyste,
czynnik $m$ daje zero w $\F_2$.
Jeśli $m$ jest nieparzyste, liczba
$m+1=2q$ jest parzysta.
Z równości $(1+z)^{2q}=(1+z^2)^q$
w $\F_2[z]$ wynika, że współczynnik
przy $z^3$ jest zerowy, czyli
$\binom{m+1}3\equiv0\pmod2$.
W obu przypadkach wynik jest więc zerowy.
To potwierdza równość obu operatorów
\emph{na} $H^*(\R P^\infty;\F_2)$,
ale nie dowodzi jej dla dowolnej przestrzeni.
\end{przyklad}

\begin{przyklad}[Niezerowe sprawdzenie relacji Adema]
\label{prz:adem-niezerowy-v8}
Dla $a=1$ i $b=2$ wzór Adema daje
$\operatorname{Sq}^1\operatorname{Sq}^2=\operatorname{Sq}^3$,
gdyż jedynym współczynnikiem jest $\binom11=1$ w $\F_2$.
Sprawdźmy tę tożsamość na $t^3\in H^3(\R P^\infty;\F_2)$.
Z formuły dwumianowej wynikają kolejne równości:
\[
 \operatorname{Sq}^2(t^3)=\binom32t^5=t^5,
 \qquad\operatorname{Sq}^1(t^5)=\binom51t^6=t^6,
 \qquad\operatorname{Sq}^3(t^3)=\binom33t^6=t^6.
\]
Obie strony są \emph{niezerowe}, co jest silniejszym testem niż
porównanie dwóch zerowych wyrażeń. Dla $a=2$, $b=3$
otrzymujemy również
$\operatorname{Sq}^2\operatorname{Sq}^3
=\operatorname{Sq}^5+\operatorname{Sq}^4\operatorname{Sq}^1$.
Na $t^3$ lewa strona jest równa
$\binom33\binom62t^8=t^8$, a prawa
$\binom35t^8+\binom31\binom44t^8=t^8$.
Te rachunki potwierdzają współczynniki w konkretnym
pierścieniu, ale nie zastępują uniwersalnego dowodu relacji.
\end{przyklad}

\begin{twierdzenie}[Postać dopuszczalna jednomianów]
Każde złożenie dodatnich kwadratów Steenroda można,
wykorzystując relacje Adema, przedstawić jako sumę
złożeń
$\operatorname{Sq}^{i_1}\cdots\operatorname{Sq}^{i_r}$
o indeksach $i_j\ge2i_{j+1}$. Takie złożenia nazywamy
\emph{dopuszczalnymi}. Twierdzenie dotyczy rozpięcia
algebry Steenroda przez jednomiany dopuszczalne;
ich liniowa niezależność wymaga odrębnego argumentu.
\end{twierdzenie}
\begin{proof}
Jeżeli $i_j<2i_{j+1}$, relacja Adema zastępuje
parę $\operatorname{Sq}^{i_j}\operatorname{Sq}^{i_{j+1}}$
sumą złożeń $\operatorname{Sq}^{i_j+i_{j+1}-k}\operatorname{Sq}^{k}$,
gdzie $k\le\lfloor i_j/2\rfloor$.
Lewa liczba jest większa od prawej co najmniej dwukrotnie,
a indeksy zachowują sumę.
Ściśle uzasadnijmy zakończenie tej procedury. Dla słowa
$\operatorname{Sq}^{i_1}\cdots\operatorname{Sq}^{i_\ell}$
o dodatnich indeksach określmy sumy końcowe
$s_j=i_j+\cdots+i_\ell$ dla $1\le j\le\ell$.
W pierwszej kolejności porównujemy długość słów, a przy
tej samej długości porównujemy leksykograficznie ciągi
$(s_\ell,s_{\ell-1},\ldots,s_1)$. W podstawieniu
pary $(a,b)$ przez $(a+b-j,j)$, gdzie $j>0$,
wszystkie sumy końcowe na prawo od tej pary pozostają
niezmienione, zaś suma zaczynająca się od drugiego
wyrazu pary maleje z $b+S$ do $j+S$, ponieważ
$j\le\lfloor a/2\rfloor<b$; $S$ jest sumą
indeksów leżących jeszcze dalej na prawo.
Każde takie podstawienie ściśle zmniejsza więc
wskazany ciąg leksykograficznie. Gdy $j=0$, otrzymujemy
$\operatorname{Sq}^{a+b}\operatorname{Sq}^0=
\operatorname{Sq}^{a+b}$, toteż długość słowa maleje.
Ponieważ całkowita suma indeksów nie zmienia się,
a długość nigdy nie rośnie, rozpatrujemy tylko skończenie
wiele możliwych słów. Nie może zatem istnieć nieskończony
ciąg redukcji.
\end{proof}

\section{Są takie przestrzenie, których pierścień kohomologii nie rozróżnia}
\begin{lemat}[Iloczyny na zawieszeniu]
\label{lem:susp-cupzero-v8}
Jeżeli $X$ jest dobrze bazowanym CW-kompleksem,
to iloczyn dwóch klas dodatnich stopni w
$\widetilde H^*(\Sigma X;R)$ jest zerowy
(dla dowolnego pierścienia współczynników $R$).
\end{lemat}
\begin{proof}
Zawieszenie $\Sigma X$ jest przestrzenią ko-H: posiada
odwzorowanie uszczypnięcia
$\Sigma X\to\Sigma X\vee\Sigma X$.
Złożenie z włączeniem bukietu do iloczynu jest
homotopijne do przekątnej
$\Sigma X\to\Sigma X\times\Sigma X$.
Gdy $a,b$ mają dodatnie stopnie, ich zewnętrzny
iloczyn ograniczony do bukietu jest zerowy,
ponieważ na każdym składniku bukietu jeden czynnik
odwzorowuje się na punkt bazowy. Ponieważ iloczyn kubkowy
jest przeciwobrazem zewnętrznego po przekątnej,
$a\smile b=0$.
\end{proof}

\begin{twierdzenie}[Zawieszona przestrzeń rzutowa nie jest bukietem sfer]
\label{thm:sigma-cp2-v8}
Niech $Y=\Sigma\C P^2$. Wtedy $Y$ i $S^3\vee S^5$
mają izomorficzne całkowite grupy homologii i izomorficzne
pierścienie kohomologii (ze współczynnikami $\Z$ i $\F_2$),
ale nie są homotopijnie równoważne.
\end{twierdzenie}
\begin{proof}
Rozkład komórkowy $\C P^2\simeq S^2\cup_\eta e^4$
daje $Y\simeq S^3\cup_{\Sigma\eta}e^5$.
Obie porównywane przestrzenie mają po jednej komórce
niezerowej w stopniach $3$ i $5$, więc z różniczek
komórkowych otrzymujemy
$\widetilde H_3=\widetilde H_5\cong\Z$ oraz pozostałe
zredukowane grupy homologii równe zeru.
Iloczyny dodatnich stopni są zerowe na obu przestrzeniach:
na $Y$ z Lematu~\ref{lem:susp-cupzero-v8}, a na bukiecie
sfer bezpośrednio.

Niech $u\in H^2(\C P^2;\F_2)$ będzie redukcją klasy
hiperpłaszczyzny, $u^2\ne0\in H^4(\C P^2;\F_2)$.
Z poprzedniego rozdziału wiemy, że
$\operatorname{Sq}^2(u)=u^2$.
Operacje Steenroda są \emph{stabilne}, to znaczy
komutują z izomorfizmem zawieszenia
$\sigma:\widetilde H^q(X;\F_2)\to
\widetilde H^{q+1}(\Sigma X;\F_2)$.
Dlatego
$\operatorname{Sq}^2(\sigma u)=\sigma(u^2)\ne0$
w $H^5(Y;\F_2)$.
Tymczasem na $S^3\vee S^5$ klasa stopnia $3$
pochodzi z $S^3$, a
$\operatorname{Sq}^2:H^3(S^3;\F_2)\to H^5(S^3;\F_2)$
jest zerowa, więc na bukiecie również.
Równoważność homotopijna zachowywałaby operacje
naturalne, zatem nie może istnieć.
\end{proof}

\begin{uwaga}[Dlaczego stabilność kwadratów jest istotna]
Stabilność $\operatorname{Sq}^i$ można sprawdzać na
unormowanym modelu wyższych diagonalnych dla par
$(CX,X)$ i na łączącej w ciągu kohomologii względnej.
Pełna zgodność łańcuchowa wymaga tego samego wyboru
spójnych diagonalnych, który pojawia się w formule Cartana.
Używamy tutaj standardowego twierdzenia o stabilności,
a nie wyprowadzamy go wyłącznie z naturalności operacji.
\end{uwaga}

\section{Wiązki wektorowe i charakterystyczny wzór Thoma}
\begin{twierdzenie}[Kwadraty Steenroda klasy Thoma]
\label{thm:thom-wu-v8}
Niech $B$ będzie parazwartym CW-kompleksem, a
$\xi\to B$ rzeczywistą wiązką wektorową rzędu $r\ge0$.
Wybierzmy metrykę włóknową i oznaczmy przez $D(\xi)$
i $S(\xi)$ odpowiednio wiązki dysków i sfer jednostkowych.
Niech $U\in H^r(D(\xi),S(\xi);\F_2)$ oznacza jej
mod-$2$ klasę Thoma. Jej ograniczenie do każdego
włókna jest generatorem grupy
$H^r(D^r,S^{r-1};\F_2)$.
Przy $r=0$ stosujemy konwencję $D^0=\{*\}$,
$S^{-1}=\varnothing$, toteż $D(\xi)=B$,
$S(\xi)=\varnothing$, a $U=1_B\in H^0(B;\F_2)$.
Dla $0\le i\le r$ istnieją jednoznaczne klasy
$w_i(\xi)\in H^i(B;\F_2)$ takie, że
$\operatorname{Sq}^i(U)=p^*w_i(\xi)\smile U$.
Są to klasy Stiefela--Whitneya; $w_0=1$,
a $w_i=0$ dla $i>r$.
\end{twierdzenie}
\begin{proof}
Współczynniki w $\F_2$ eliminują problem
orientacji włókna: każdy automorfizm
$H^r(D^r,S^{r-1};\F_2)\cong\F_2$ jest
tożsamością. W lokalnych trywializacjach
klasy generatorów sklejają się więc w
mod-$2$ klasę Thoma także dla wiązki
nieorientowalnej. Względny izomorfizm Thoma
$H^i(B;\F_2)\xrightarrow{\cong}
H^{r+i}(D(\xi),S(\xi);\F_2)$,
$x\mapsto p^*x\smile U$,
pozwala jednoznacznie zdefiniować $w_i$ jako przeciwobraz
$\operatorname{Sq}^i(U)$. Z własności
$\operatorname{Sq}^0=\id$ dostajemy $w_0=1$.
Z niestabilności operacji w kohomologii względnej
$\operatorname{Sq}^i(U)=0$ dla $i>r$,
a więc $w_i=0$ dla $i>r$. Można tu
utożsamić $H^*(D(\xi),S(\xi);\F_2)$ z
kohomologią zredukowaną ilorazu
$D(\xi)/S(\xi)$, jeśli $r\ge1$,
i zastosować niestabilność do klasy stopnia $r$.
Dla $r=0$ ta sama równość wynika z
$\operatorname{Sq}^i(1_B)=0$ dla $i>0$.
Naturalność klas wynika z naturalności izomorfizmu
Thoma i operacji Steenroda.
\end{proof}

\begin{wniosek}[Wzór Whitney'a]
Dla dwóch wiązek rzeczywistych $\xi,\eta\to B$
zachodzi $w(\xi\oplus\eta)=w(\xi)\smile w(\eta)$,
gdzie $w=\sum_iw_i$.
\end{wniosek}
\begin{proof}
Ustalmy rzędy $r$ i $s$ wiązek $\xi$ i $\eta$.
Dla iloczynu włóknistego $D(\xi)\times_B D(\eta)$
rozpatrujemy względny iloczyn klas
$U_\xi$ i $U_\eta$, zdefiniowany przez
iloczyn zewnętrzny oraz cofnięcie na iloczyn
włóknisty. Po utożsamieniu pary włóknistej
z odpowiednią parą dysków wiązki
$\xi\oplus\eta$ otrzymujemy klasę
$U_{\xi\oplus\eta}$: na każdej parze włókien
jej ograniczenie jest iloczynem dodatnich
generatorów względnych grup w stopniach
$r$ i $s$, a zatem generatorem w stopniu
$r+s$. Wspomniane utożsamienie można
otrzymać przez ciągłą deformację produktu
dysków na dysk jednostkowy w sumie prostej;
obejmuje ono również rząd zero dzięki
przyjętej konwencji.

Naturalność kwadratów Steenroda i zewnętrzna
postać wzoru Cartana dają, w kohomologii
tej pary względnej,
\[
\begin{aligned}
\operatorname{Sq}(U_{\xi\oplus\eta})
 &=\operatorname{Sq}(U_\xi)\smile
   \operatorname{Sq}(U_\eta)\\
 &=(p^*w(\xi)\smile U_\xi)\smile
   (p^*w(\eta)\smile U_\eta)\\
 &=p^*(w(\xi)\smile w(\eta))
   \smile U_{\xi\oplus\eta}.
\end{aligned}
\]
Powyższe mnożenia rozumiemy w odpowiednich
parach względnych, po wspomnianych cofnięciach;
nad $\F_2$ nie występują dodatkowe znaki.
Z definicji $w(\xi\oplus\eta)$ lewa strona
jest także równa
$p^*w(\xi\oplus\eta)\smile U_{\xi\oplus\eta}$.
Różnowartościowość izomorfizmu Thoma dla
$\xi\oplus\eta$ daje żądany wzór Whitney'a.
\end{proof}

\begin{przyklad}[Wiązka tautologiczna nad $\R P^n$]
Dla rzeczywistej wiązki liniowej tautologicznej
$\gamma^1\to\R P^n$ z generatorem
$t\in H^1(\R P^n;\F_2)$ mamy $w(\gamma^1)=1+t$.
Ponieważ $\gamma^1\oplus(\gamma^1)^\perp$ jest trywialna,
ze wzoru Whitney'a otrzymujemy
$w((\gamma^1)^\perp)=(1+t)^{-1}=1+t+\cdots+t^n$.
Równości zachodzą w pierścieniu
$\F_2[t]/(t^{n+1})$; nie zakładamy zbieżności
nieskończonego szeregu formalnego.
\end{przyklad}

\section{Co wymaga dalszej techniki?}
Sam wzór Cartana jest w pełni obliczalny na klasach
pierścieni generowanych przez elementy małego stopnia;
relacje Adema pozwalają upraszczać iterowane kwadraty.
Pełny niezależny dowód ogólnych relacji Adema wymaga
budowy spójnych działań symetrycznych na iterowanych
produktach i zostaje wyraźnie oddzielony od podanych
rachunków. W następnym rozdziale stabilność operacji
będzie kluczem do niezerowości zawieszonego odwzorowania Hopfa.
